% GENERATED FILE. Edit canonical lessons, then run npm run generate:books.
% canonical-source-hash: fnv1a64:7c4d6ca09f275073
% canonical-lessons: SA-C176-parivartayati, SA-C176-nivedayati, SA-C176-amantrayati, SA-C176-sodhayati, SA-C176-purayati

\chapter{Changes, Reports, Invites, Cleans, Fills}
\label{ch:sa-changes-reports-invites-cleans-fills}
\hlchaptermodality{176}

\begin{chapteropening}
\textbf{By the end of this chapter:} \emph{I can say changes, reports, invites, cleans and fills.}

Complete the last lesson of chapter 176: I can say changes, reports, invites, cleans and fills.
\end{chapteropening}

\section[\texorpdfstring{\sk{परिवर्तयति} (parivartayati)}{परिवर्तयति (parivartayati)}]{\sk{परिवर्तयति} (parivartayati) — changes}

\label{lesson:SA-C176-parivartayati}



\begin{quote}
Before the new one: say the Sanskrit for endures, then the Sanskrit for tries.
\end{quote}

\subsection*{You'll want to know: \sk{परिवर्तयति}}
\textbf{\sk{परिवर्तयति}} — \emph{parivartayati} — \textquotedblleft{}changes\textquotedblright{}.

\begin{grammarlens}[title={What you've built}]
One more word for this chapter.
\end{grammarlens}

\subsection*{Your turn}
\begin{itemize}
\raggedright
  \item \emph{Say it:} \emph{parivartayati}
  \item \emph{Say it:} \emph{parivartayati}, once more
  \item \emph{Say it:} say \emph{parivartayati}
  \item \emph{Recall:} say the Sanskrit for endures, then the Sanskrit for tries, then say \emph{parivartayati} again
\end{itemize}

\begin{tcolorbox}[breakable,colback=teal!4,colframe=teal!35!black,title={Before you move on}]
What does \sk{परिवर्तयति} mean? (\textquotedblleft{}Changes\textquotedblright{}.) Say it once more.
\end{tcolorbox}
\section[\texorpdfstring{\sk{निवेदयति} (nivedayati)}{निवेदयति (nivedayati)}]{\sk{निवेदयति} (nivedayati) — reports, submits}

\label{lesson:SA-C176-nivedayati}



\begin{quote}
Before the new one: say the Sanskrit for tries, then the Sanskrit for changes.
\end{quote}

\subsection*{You'll want to know: \sk{निवेदयति}}
\textbf{\sk{निवेदयति}} — \emph{nivedayati} — \textquotedblleft{}reports, submits\textquotedblright{}.

\begin{grammarlens}[title={What you've built}]
One more word for this chapter.
\end{grammarlens}

\subsection*{Your turn}
\begin{itemize}
\raggedright
  \item \emph{Say it:} \emph{nivedayati}
  \item \emph{Say it:} \emph{nivedayati}, once more
  \item \emph{Say it:} say \emph{nivedayati}
  \item \emph{Recall:} say the Sanskrit for tries, then the Sanskrit for changes, then say \emph{nivedayati} again
\end{itemize}

\begin{tcolorbox}[breakable,colback=teal!4,colframe=teal!35!black,title={Before you move on}]
What does \sk{निवेदयति} mean? (\textquotedblleft{}Reports, submits\textquotedblright{}.) Say it once more.
\end{tcolorbox}
\section[\texorpdfstring{\sk{आमन्त्रयति} (āmantrayati)}{आमन्त्रयति (āmantrayati)}]{\sk{आमन्त्रयति} (āmantrayati) — invites}

\label{lesson:SA-C176-amantrayati}



\begin{quote}
Before the new one: say the Sanskrit for changes, then the Sanskrit for reports.
\end{quote}

\subsection*{You'll want to know: \sk{आमन्त्रयति}}
\textbf{\sk{आमन्त्रयति}} — \emph{āmantrayati} — \textquotedblleft{}invites\textquotedblright{}.

\begin{grammarlens}[title={What you've built}]
One more word for this chapter.
\end{grammarlens}

\subsection*{Your turn}
\begin{itemize}
\raggedright
  \item \emph{Say it:} \emph{āmantrayati}
  \item \emph{Say it:} \emph{āmantrayati}, once more
  \item \emph{Say it:} say \emph{āmantrayati}
  \item \emph{Recall:} say the Sanskrit for changes, then the Sanskrit for reports, then say \emph{āmantrayati} again
\end{itemize}

\begin{tcolorbox}[breakable,colback=teal!4,colframe=teal!35!black,title={Before you move on}]
What does \sk{आमन्त्रयति} mean? (\textquotedblleft{}Invites\textquotedblright{}.) Say it once more.
\end{tcolorbox}
\section[\texorpdfstring{\sk{शोधयति} (śodhayati)}{शोधयति (śodhayati)}]{\sk{शोधयति} (śodhayati) — cleans, corrects}

\label{lesson:SA-C176-sodhayati}



\begin{quote}
Before the new one: say the Sanskrit for reports, then the Sanskrit for invites.
\end{quote}

\subsection*{You'll want to know: \sk{शोधयति}}
\textbf{\sk{शोधयति}} — \emph{śodhayati} — \textquotedblleft{}cleans, corrects\textquotedblright{}.

\begin{grammarlens}[title={What you've built}]
One more word for this chapter.
\end{grammarlens}

\subsection*{Your turn}
\begin{itemize}
\raggedright
  \item \emph{Say it:} \emph{śodhayati}
  \item \emph{Say it:} \emph{śodhayati}, once more
  \item \emph{Say it:} say \emph{śodhayati}
  \item \emph{Recall:} say the Sanskrit for reports, then the Sanskrit for invites, then say \emph{śodhayati} again
\end{itemize}

\begin{tcolorbox}[breakable,colback=teal!4,colframe=teal!35!black,title={Before you move on}]
What does \sk{शोधयति} mean? (\textquotedblleft{}Cleans, corrects\textquotedblright{}.) Say it once more.
\end{tcolorbox}
\section[\texorpdfstring{\sk{पूरयति} (pūrayati)}{पूरयति (pūrayati)}]{\sk{पूरयति} (pūrayati) — fills}

\label{lesson:SA-C176-purayati}



\begin{quote}
Before the new one: say the Sanskrit for invites, then the Sanskrit for cleans.
\end{quote}

\subsection*{You'll want to know: \sk{पूरयति}}
\textbf{\sk{पूरयति}} — \emph{pūrayati} — \textquotedblleft{}fills\textquotedblright{}.

\begin{grammarlens}[title={What you've built}]
One more word for this chapter.
\end{grammarlens}

\subsection*{Your turn}
\begin{itemize}
\raggedright
  \item \emph{Say it:} \emph{pūrayati}
  \item \emph{Say it:} \emph{pūrayati}, once more
  \item \emph{Say it:} say \emph{pūrayati}
  \item \emph{Recall:} say the Sanskrit for invites, then the Sanskrit for cleans, then say \emph{pūrayati} again
\end{itemize}

\begin{tcolorbox}[breakable,colback=teal!4,colframe=teal!35!black,title={Before you move on}]
What does \sk{पूरयति} mean? (\textquotedblleft{}Fills\textquotedblright{}.) Say it once more.
\end{tcolorbox}
